\documentclass[conference]{IEEEtran}
\IEEEoverridecommandlockouts
% The preceding line is only needed to identify funding in the first footnote. If that is unneeded, please comment it out.
\usepackage{cite}
\usepackage{amsmath,amssymb,amsfonts}
\usepackage{algorithmic}
\usepackage{graphicx}
\usepackage{textcomp}
\usepackage{xcolor}
\usepackage{booktabs}
\usepackage{enumitem}
\newcommand{\Zhaozhuo}[1]{\textcolor{red}{[Zhaozhuo: #1]}}
\def\BibTeX{{\rm B\kern-.05em{\sc i\kern-.025em b}\kern-.08em
    T\kern-.1667em\lower.7ex\hbox{E}\kern-.125emX}}
\begin{document}

\title{ICDM Workshop: VISTA: Visionary Innovation in Standards and Technology of GenAI
}

\maketitle

\begin{abstract}
The VISTA: Visionary Innovation in Standards and Technology of GenAI workshop at ICDM 2025 examines the intersection of Generative AI (GenAI) advancements and evolving community and legal standards. As GenAI transforms industries, it raises challenges in legal compliance and common sense considerations. This workshop convenes researchers, industry experts, and legal scholars to explore the alignment between technical innovation and regulatory frameworks, fostering discussions on responsible GenAI development.
\end{abstract}

\begin{IEEEkeywords}
GenAI, technology, standards.
\end{IEEEkeywords}

\section{Duration}
Half-day workshop.

\section{Topic Description}
Generative AI (GenAI) models, including large language models (LLMs) and diffusion models, are reshaping industries by automating content creation, enhancing decision-making, and enabling new forms of user interaction. However, their widespread adoption presents challenges in both technology and standards. How these new technologies perform with legal compliance and respect to common sense remains to be explored. This workshop brings together researchers and practitioners to explore the interaction of advanced GenAI technology and evolving corresponding standards. It aims to foster discussions on future innovations of GenAI in a regulated landscape.

\noindent\textbf{Topic of Interest.} We welcome submissions on innovations in GenAI at the intersection of technology and standards. We are particularly interested in contributions that align GenAI technologies with legal regulations and common sense considerations.
Potential topics include but are not limited to:

\begin{itemize}[leftmargin=*]
\item New GenAI models with legal considerations
\item Novel GenAI model training and inference paradigm with legal compliance
\item Efficient GenAI model customization with respect to legal regulations and common sense
\item Analysis of data influence on GenAI models' behaviors related to common sense
\item Copyright challenges in GenAI, including detection, quantification, and mitigation techniques
\item Machine unlearning of GenAI models for legal compliance
\item Implementations of GenAI standards in various applications
\item Case studies of legal challenges in GenAI applications.
\item GenAI-powered applications in digital finance
\item Technical contributions to ensure GenAI's compliance with financial regulations (e.g., GDPR, SEC requirements)
\item GenAI applications in legal research, including contract analysis and case prediction
\item GenAI applications in risk management and regulatory monitoring
\end{itemize}

This workshop aims to bring together technical researchers, legal experts, and financial regulators to collaborate on developing comprehensive and practical approaches to GenAI development and its applications. The focus will be on ensuring legal compliance while integrating common-sense considerations.

\section{Submission Details}
We invite both research and vision papers to contribute to the intersection of technology and standards in GenAI. Submissions should follow the ICDM regular research paper format guidelines and be submitted via OpenReview.
Accepted papers will be presented at the workshop in either oral or poster presentations. The reviewing process will be double-blind. No formal proceedings will be generated from this workshop. Accepted papers will be published on OpenReview.
\begin{itemize}
    \item \textbf{Paper format:} Up to 6 pages (excluding references and appendices)
    \item \textbf{Submission deadline:} September 5th, 2025
    \item \textbf{Notification of acceptance:} October 15th, 2025
    \item \textbf{Camera-ready deadline:} November 1st, 2025
    \item \textbf{Workshop date:} November 12th, 2025 (tentative)
\end{itemize}
    
\section{Plans to Attract Quality Submissions}
Besides outreach and dissemination, the organizers plan to attract high-quality submissions through the following efforts.

\noindent \textbf{Invited Papers.} We plan to invite leading researchers in data mining, AI, and social sciences to submit their work to the workshop. These research groups have informally agreed to contribute to our workshop.
\begin{itemize}
    \item Prof. Xia Ben Hu at Rice University
    \item Prof. Anshumali Shrivastava at Rice University
    \item Prof. Yanjie Fu at Arizona State University
    \item Prof. Kunpeng Liu at Portland State University
    \item Prof. Kaixiong Zhou at North Carolina State University
    \item Prof. Ruixiang Tang at Rutgers University
\end{itemize}

\noindent \textbf{Journal Tracks.}  We plan to invite authors published in data mining journals, such as ACM Transactions on Knowledge Discovery from Data, IEEE Transactions on Knowledge and Data Engineering, and IEEE Transactions on Big Data, to present their research on the interdisciplinary research of technology and standards of GenAI.

\noindent \textbf{Industry Demonstrations.} Our organizers plan to invite industry collaborators, including the Fin AI, Amazon, and Google, to present their research and implementation on novel GenAI applications, including both law service and digital finance.

\noindent \textbf{Closed-Door Roundtable Discussion and Survey Drafting.} We will host a closed-door roundtable discussion to explore future challenges and opportunities in GenAI technology and standards. This session will contribute to drafting a survey paper that outlines key insights and directions for the field. Authors of accepted papers will be invited to participate.

\section{Preliminary List of Invited Speakers}
Preliminary list of invited speakers include:
\begin{itemize}
    \item \textbf{Opening Remarks \& Keynote Speaker:} Jing Gao, Associate Professor at Purdue University
    \item \textbf{Tentative Keynote Speaker, GenAI Technology:} Vladimir Braverman, Professor at Johns Hopkins University
    \item \textbf{Tentative Keynote Speaker, GenAI Standards:} David Atkinson, J.D., Lecturer at the McCombs School of Business at The University of Texas at Austin
    \item \textbf{Industry Showcase Speaker:} Jimin Huang from the Fin AI
\end{itemize}

\section{Biographies of Organizers }
\noindent \textbf{Denghui Zhang} is an Assistant Professor in the School of Business at Stevens Institute of Technology. He studies the interplay between LLMs/GenAI, legal and socioethical issue, business and innovation. His work is published in refereed venues, such as IEEE TKDE, SIGKDD, NeurIPS, EMNLP, AAAI, etc., and won Best Student Paper award at 2023 International Conference on Information Systems.

\noindent \textbf{Zhaozhuo Xu} is an assistant professor at Stevens Institute of Technology, focusing on scalable and sustainable ML. His work, published in venues like NeurIPS, ICML, and ICLR, has been adopted by Huggingface and startups.  He also serves as an area chair for major conferences, including ICLR, ICML, ARR, EMNLP and COLING. He received the AAAI-25 New Faculty Highlights.

\noindent \textbf{Jimin Huang} is the president of The Fin AI community, which advances open science, tooling, and model development for financial services with a focus on responsible innovation. His research spans NLP and computational finance, emphasizing financial LLMs and open-source contributions. He organizes the FinLLM Challenge at FinNLP-AgentScen @ IJCAI-2024 and serves as general chair for the Joint Workshop on FinNLP, FNP, and LLMFinLegal @ COLING 2025.

\noindent \textbf{Qingyun Wang} is the incoming assistant professor at William \& Mary. He is among the first researchers to develop a virtual scientific research assistant for literature-based discovery by extracting and synthesizing insights from papers. He received the NAACL-HLT 2021 Best Demo Reward. He has experience presenting tutorials at EMNLP 2021, EACL 2024, and LREC-COLING 2024. He organized AI4Research workshop at AAAI 2025 and the Language + Molecules Workshop at ACL 2024. 

\noindent \textbf{Jing Gao} is an Associate Professor at Purdue University's Elmore Family School of Electrical and Computer Engineering. Her research spans data mining, focusing on information veracity, crowdsourcing, knowledge graphs, anomaly detection, and multi-source data integration, with applications in healthcare, cybersecurity, and education.
She has received the NSF CAREER Award, ICDM Tao Li Award, and SDM/IBM Early Career Data Mining Researcher Award. She also serves as an editor for ACM Transactions on Intelligent Systems and Technology and IEEE Transactions on Knowledge and Data Engineering.

\section{Organizing Committee}
\begin{itemize}
    \item \textbf{General Chairs:} Jing Gao and Denghui Zhang
    \item \textbf{Program Chair:} Zhaozhuo Xu
    \item \textbf{Local Chair:} Qingyun Wang
    \item \textbf{Publicity Chair:} Jimin Huang
\end{itemize}

\section{Preliminary Program Committee}
The program committee includes:
\begin{itemize}
    \item David Atkinson (UT Austin)
    \item Qianqian Xie (the Fin AI)
    \item Boyi Wei (Princeton)
    \item Tianfu Fu (OpenAI)
    \item Jialiang Xu (Stanford)
    \item Yanzhou Pan (Google)
    \item Jiateng Liu (UIUC)
\end{itemize}

\section{Contact Information of the Organizers}
\begin{itemize}
    \item Denghui Zhang, Assistant Professor, School of Business, Stevens Institute of Technology. Email: dzhang42@stevens.edu
    \item  Zhaozhuo Xu, Assistant Professor, Department of Computer Science, Stevens Institute of Technology. Email: zxu79@stevens.edu
    \item Jimin Huang, President of the Fin AI. Email: jimin.huang@thefin.ai
    \item Qingyun Wang, Incoming Assistant Professor, College of William \& Mary. Email: qingyun4@illinois.edu
    \item Jing Gao, Associate Professor at Purdue University. Email: jinggao@purdue.edu
\end{itemize}

\section{Past Workshop Organizing Experiences}
Jing Gao is one of the leading organizers of the TrueFact workshop held in conjunction with ACM KDD from 2019 to 2022. The workshop provided a forum for researchers and practitioners in academia, government and industry to share insights on topics related to information quality. 

Zhaozhuo Xu is the leading organizer for the Research On Algorithms \& Data Structures (ROADS) to Mega-AI Models Workshop at MLSys 2023\footnote{https://roads2megaai.github.io/main.html}, which invites:
\begin{itemize}
    \item Jonathan Frankle, Chief Scientist, Databricks Mosaic AI
    \item Michael Mitzenmacher, Professor, Harvard University
    \item Furong Huang, Associate Professor, University of Maryland, College Park
    \item Bilge Acun, Research Scientist at Meta AI / FAIR
\end{itemize}
Moreover, the ROADs workshop includes oral paper presentations from CMU, Stanford, Airbnb, and Amazon. The total attendees rank top in the MLSys 2023 workshops.

\section{Schedule}

\begin{table}[h!]
\renewcommand{\arraystretch}{1.1}
\setlength{\tabcolsep}{6pt}
\centering
\caption{Proposed Workshop Schedule}
\begin{tabular}{@{}ll@{}}
\toprule
\textbf{Time}       & \textbf{Event} \\
\midrule
9:30--9:40   & Opening Remarks \\
9:40--10:20  & \textit{Keynote:} Jing Gao, Purdue University \\
10:20--11:00 & \textit{Keynote:} Vladimir Braverman, Johns Hopkins University \\
11:00--11:40 & \textit{Keynote:} David Atkinson, J.D., UT Austin \\
11:40--12:10 & \textit{Keynote:} Ang Li, University of Maryland \\
12:10--13:30 & Lunch Break \\
13:30--14:00 & \textit{Industry Showcase:} Jimin Huang, The Fin AI \\
14:00--14:30 & Roundtable: Perspective White Paper Discussion \\
14:30--16:00 & Poster Session \\
\bottomrule
\end{tabular}
\end{table}

\end{document} 